\documentclass[10pt]{article}

\usepackage[a4paper,margin=1.8cm]{geometry}
\usepackage{amsmath,amssymb}
\usepackage{booktabs}
\usepackage{array}
\usepackage{graphicx}
\usepackage{microtype}
\usepackage[hidelinks]{hyperref}
\usepackage{enumitem}
\usepackage{xcolor}
\usepackage{longtable}
\usepackage{placeins}

\hypersetup{
  pdftitle={Supplementary Information: A Paired, Floor-Guarded Evaluation Protocol for INT8 and FP8 Object Detectors under Image Corruptions},
  pdfauthor={Dinh Thuan Nguyen; Lam Phuong Nguyen; Vinh Huy Nguyen; Sy Vu Quang; Mohan Rajesh Elara; Anh Vu Le},
  pdfsubject={Supplementary methods, evidence summaries, and reproducibility contract}
}

\newcommand{\AP}{\mathrm{AP}}
\newcommand{\Jclean}{\mathrm{J95}}
\newcommand{\E}{\mathcal{E}}
\newcommand{\DeltaE}{\Delta E}
\newcommand{\DeltaR}{\Delta R}
\newcommand{\DeltaPsi}{\Delta\Psi}
\renewcommand{\thesection}{S\arabic{section}}
\renewcommand{\thetable}{S\arabic{table}}
\renewcommand{\thefigure}{S\arabic{figure}}
\renewcommand{\theequation}{S\arabic{equation}}
\setlist[itemize]{leftmargin=*,nosep}
\setlength{\tabcolsep}{4pt}

\title{Supplementary Information for\\[1mm]
\textbf{A Paired, Floor-Guarded Evaluation Protocol for INT8 and FP8 Object Detectors under Image Corruptions}}
\author{Dinh Thuan Nguyen$^{1}$, Lam Phuong Nguyen$^{1}$, Vinh Huy Nguyen$^{2}$,\\
Sy Vu Quang$^{4}$, Mohan Rajesh Elara$^{3}$, Anh Vu Le$^{4,*}$\\[2mm]
\small $^1$Faculty of Electrical and Electronics Engineering, Ton Duc Thang University, Vietnam\\
\small $^2$Center of Visualization and Simulation, Duy Tan University, Vietnam\\
\small $^3$ROAR Lab, Singapore University of Technology and Design, Singapore\\
\small $^4$Advanced Intelligent Technology Research Group, Ton Duc Thang University, Vietnam\\
\small $^*$Corresponding author: leanhvu@tdtu.edu.vn}
\date{}

\begin{document}
\maketitle

\section{Purpose, scope, and evidence hierarchy}

This supplement documents the estimands, diagnostic summaries, sensitivity
analyses, and artifact contract supporting the main article. It does not add a
new population-level claim. The evidence hierarchy used throughout the paper is:
(i) the 144-cell four-dataset YOLO landscape for exploratory factorial
characterization; (ii) the selection-disjoint VOC/KITTI final holdout as the strongest
split-independent conditional evidence; (iii) targeted realization, training,
calibration, class-universe, codec, weighting, and runtime sensitivities; and
(iv) RT-DETR and RetinaNet experiments as recipe-portability stress cases.

The labels \emph{INT8} and \emph{FP8} denote the complete recorded TensorRT
treatments---quantizer placement, Q/DQ coverage, calibration or scale
construction, graph rewriting, mixed-precision fallbacks, engine building, and
frozen decoding---rather than datatype effects in isolation. All quantities are
reported in AP points unless stated otherwise.

\section{Atomic paired design and estimands}

For a fixed dataset $d$, checkpoint $m$, endpoint $b$, and corruption condition
$c$, let $\Jclean$ denote the deterministic JPEG-quality-95 matched-clean control. The explicit $\Jclean$ label avoids confusion with AP at IoU 0.95.
For quantized treatment $q$, define the FP32-referenced quantization gap and the
clean-adjusted excess gap as
\begin{align}
Q_{d,m,q,b,c}
 &= \AP_{d,m,\mathrm{FP32},b,c}-\AP_{d,m,q,b,c},\\
E_{d,m,q,b,c}
 &= Q_{d,m,q,b,c}-Q_{d,m,q,b,\Jclean}.
\end{align}
A positive $E$ means that corruption widens the treatment--FP32 discrepancy
relative to matched clean; a negative $E$ means that the discrepancy contracts.
Neither sign is, by itself, evidence of retained absolute accuracy.

The direct treatment interaction is computed inside each common-image resample:
\begin{align}
\DeltaE_{d,m,b,c}
 &= E_{d,m,\mathrm{INT8},b,c}-E_{d,m,\mathrm{FP8},b,c}\\
 &= \bigl(\AP_{\mathrm{FP8},c}-\AP_{\mathrm{INT8},c}\bigr)
  -\bigl(\AP_{\mathrm{FP8},\Jclean}-\AP_{\mathrm{INT8},\Jclean}\bigr).
\end{align}
The FP32 term cancels algebraically. Consequently, $\DeltaE$ asks whether the
FP8--INT8 discrepancy changes under corruption after removing the discrepancy
already present on matched clean inputs. A positive value indicates expansion;
a negative value indicates contraction. This is a descriptive two-by-two
interaction contrast, not a causal difference-in-differences estimator.

For area endpoints, the within-treatment size contrast is
\begin{equation}
\Psi_{d,m,q,c}=E_{d,m,q,S,c}-E_{d,m,q,L,c},
\end{equation}
with direct size interaction $\DeltaPsi=\Psi_{\mathrm{INT8}}-
\Psi_{\mathrm{FP8}}$. TT100K uses its native height-based small-like and
large-like endpoints, which are reported separately and never pooled
numerically with the COCO/VOC/KITTI area endpoints.

The signed clean-to-corrupted AP loss used by the absolute-accuracy guardrail is
\begin{equation}
D_{d,m,p,b,c}=\AP_{d,m,p,b,\Jclean}-\AP_{d,m,p,b,c}.
\end{equation}
Thus $D>0$ is an AP loss and $D<0$ is an AP gain; despite its role in the
absolute-accuracy guardrail, $D$ is not an absolute value.
The direct interaction can therefore also be written as
\begin{equation}
\Delta E_{d,m,b,c}=D_{d,m,\mathrm{INT8},b,c}-D_{d,m,\mathrm{FP8},b,c}.
\end{equation}
The protocol is therefore \emph{floor-guarded}: $D$, corrupted AP, and low-AP
inventories are reported beside relative interactions. It is not a formal
statistical correction for censoring or an AP floor.

As a secondary scale sensitivity, relative retention normalizes each corrupted
arm by that treatment's own matched-clean AP:
\begin{align}
R_{d,m,q,b,c}
 &=100\frac{\AP_{d,m,q,b,c}}{\AP_{d,m,q,b,\Jclean}},\\
\DeltaR_{d,m,b,c}
 &=R_{d,m,\mathrm{FP8},b,c}-R_{d,m,\mathrm{INT8},b,c}.
\end{align}
$\DeltaR$ is reported in percentage points (pp); positive values mean that FP8
retains a larger percentage of its own matched-clean AP. Because this ratio
changes the scale and effective weighting induced by clean AP, it is a
sensitivity estimand rather than a replacement for the AP-point interaction
$\DeltaE$.

The image is the resampling unit. Each direct comparison reuses one deterministic
schedule of 2,000 image-index bootstrap draws across all four INT8/FP8 by
clean/corrupted arms. Percentile intervals quantify finite-evaluation-set
uncertainty conditional on the frozen datasets, checkpoints, calibration lists,
engines, corruption materializations, and decoder. They do not propagate a
population distribution over datasets, training runs, calibration runs, or
corruption generators.

\section{Decision-impact diagnostic}

The methodological value of clean adjustment can be measured directly. For each
of the 144 exploratory cells, the raw corrupted gap satisfies
\begin{equation}
(\AP_{\mathrm{FP8}}-\AP_{\mathrm{INT8}})_{c}
=(\AP_{\mathrm{FP8}}-\AP_{\mathrm{INT8}})_{\Jclean}+\DeltaE.
\end{equation}
The raw corrupted comparison favored FP8 in 134 cells, but 56 of those cells
had negative $\DeltaE$ point estimates after subtracting the matched-clean
discrepancy (Table~\ref{tab:decision-impact}). Of these 56 point-sign
reversals, 17 paired percentile intervals lay entirely below zero and 39 crossed
zero; 56 is therefore not a count of statistically resolved discoveries. Thus
a raw corrupted leaderboard comparison
would conflate clean fidelity with corruption-specific interaction in 38.9\% of
the finite grid.

\begin{table}[htbp]
\centering
\caption{Decision-impact inventory over the 144 exploratory cells. ``Below AP''
counts a cell when at least one quantized treatment crosses the stated floor.}
\label{tab:decision-impact}
\small
% Generated by scripts/make_decision_impact.py; do not edit manually.
\begin{tabular}{lr}
\toprule
Diagnostic & Cells \\
\midrule
Raw corrupted FP8--INT8 gap $>0$ & 134/144 \\
Raw corrupted FP8--INT8 gap $<0$ & 10/144 \\
Clean-adjusted $\Delta E>0$ & 78/144 \\
Clean-adjusted $\Delta E<0$ & 66/144 \\
Raw-positive but adjusted-negative & 56/144 \\
At least one format below 10 AP & 18/144 \\
At least one format below 5 AP & 5/144 \\
\bottomrule
\end{tabular}

\end{table}

\begin{table}[htbp]
\centering
\caption{Sensitivity of the 56 positive-raw-gap to negative-$\DeltaE$
point-sign reversals to minimum interaction magnitude. Intervals are paired
image-bootstrap percentile intervals and are not multiplicity-adjusted.}
\label{tab:decision-reversal-sensitivity}
\small
\resizebox{0.86\textwidth}{!}{% Generated by analysis/build_cviu_revision_artifacts.py.
\begin{tabular}{lrrr}
\toprule
Point-effect filter (AP) & Reversals & Interval below zero & Interval crosses zero \\
\midrule
$|\Delta E|\geq 0.00$ & 56 & 17 & 39 \\
$|\Delta E|\geq 0.25$ & 34 & 17 & 17 \\
$|\Delta E|\geq 0.50$ & 20 & 14 & 6 \\
$|\Delta E|\geq 1.00$ & 11 & 10 & 1 \\
\bottomrule
\end{tabular}
}
\end{table}

\FloatBarrier

\section{Quantized graph coverage}

Table~\ref{tab:graph-coverage} summarizes hash-bound ONNX graph coverage. These
counts characterize inserted Q/DQ structure and zero-point datatype; they are
not a claim that every TensorRT kernel executes at the nominal precision. That
stronger claim would require per-layer engine inspection or profiling, which is
outside the recorded evidence package.

\begin{table}[htbp]
\centering
\caption{Graph-level Q/DQ coverage across the evaluated detector families.
Ranges reflect the three dataset-specific graphs for each row.}
\label{tab:graph-coverage}
\scriptsize
\resizebox{\textwidth}{!}{% Generated from hash-bound ONNX graphs; graph-level coverage, not TensorRT kernel precision.
\begin{tabular}{lllrrrrrl}
\toprule
Family & Model & Format & Graphs & Q & Weight Q & Act. Q & Direct Q/DQ nodes & Zero point \\
\midrule
YOLO11 & YOLO11m & INT8 & 3 & 437--441 & 113 & 324--328 & 352--355/405 & INT8 \\
YOLO11 & YOLO11m & FP8 & 3 & 213--219 & 101--104 & 112--115 & 133--136/405 & FLOAT8E4M3FN \\
YOLO11 & YOLO11n & INT8 & 3 & 325--341 & 85--88 & 240--253 & 263--275/320 & INT8 \\
YOLO11 & YOLO11n & FP8 & 3 & 152--158 & 69--72 & 83--86 & 99--102/320 & FLOAT8E4M3FN \\
YOLO11 & YOLO11x & INT8 & 3 & 682--690 & 174 & 508--516 & 546--552/617 & INT8 \\
YOLO11 & YOLO11x & FP8 & 3 & 346--352 & 161--164 & 185--188 & 215--218/617 & FLOAT8E4M3FN \\
RT-DETR & RT-DETR-L & INT8 & 3 & 534--538 & 197--199 & 335--339 & 351--353/963 & INT8 \\
RT-DETR & RT-DETR-L & FP8 & 3 & 398--401 & 167--169 & 231--232 & 251--253/963 & FLOAT8E4M3FN \\
RetinaNet & R50-FPN-v2 & INT8 & 3 & 217 & 111 & 106 & 139/1076 & INT8 \\
RetinaNet & R50-FPN-v2 & FP8 & 3 & 214 & 110 & 104 & 137/1076 & FLOAT8E4M3FN \\
\bottomrule
\end{tabular}
}
\end{table}

\FloatBarrier

\section{Exploratory interaction summaries}

Table~\ref{tab:dataset-summary} reports the four dataset margins, while
Table~\ref{tab:factor-summary} reports capacity, corruption, and ordinal-severity
margins. Cell-wise percentile signs are descriptive inventories rather than
multiplicity-adjusted discoveries. The balanced macro is the arithmetic mean of
the finite, defined 144-cell grid; it is not a deployment-population
weighted effect.

\begin{table}[htbp]
\centering
\caption{Dataset-level direct interaction summaries. TT100K's $\DeltaPsi$ uses
the height endpoint; the other rows use the area endpoint.}
\label{tab:dataset-summary}
\small
\resizebox{0.88\textwidth}{!}{\begin{tabular}{lrrrr}
\toprule
Dataset & Cells & Mean $\Delta E$ (AP pt) & 95\% percentile sign (+/-/cross) & Mean $\Delta\Psi$ (AP pt) \\
\midrule
COCO & 36 & -0.19 & 6 / 10 / 20 & -0.12 \\
VOC & 36 & -0.04 & 12 / 5 / 19 & -0.44 \\
KITTI & 36 & -0.46 & 5 / 11 / 20 & -0.95 \\
TT100K & 36 & +0.61 & 2 / 0 / 34 & -1.38 \\
\bottomrule
\end{tabular}
}
\end{table}

\begin{table}[htbp]
\centering
\caption{Factor-level heterogeneity in the exploratory direct ledger.}
\label{tab:factor-summary}
\scriptsize
\resizebox{\textwidth}{!}{\begin{tabular}{llrrrr}
\toprule
Factor & Level & Cells & Mean $\Delta E$ & Range $\Delta E$ & 95\% percentile sign (+/-/cross) \\
\midrule
Checkpoint rung & YOLO11n & 48 & +0.07 & -2.71 to +3.10 & 11 / 7 / 30 \\
Checkpoint rung & YOLO11m & 48 & -0.11 & -4.28 to +2.82 & 9 / 10 / 29 \\
Checkpoint rung & YOLO11x & 48 & -0.02 & -4.23 to +2.09 & 5 / 9 / 34 \\
Corruption & Gaussian noise & 36 & -0.03 & -3.66 to +2.82 & 9 / 11 / 16 \\
Corruption & Motion blur & 36 & -0.24 & -4.28 to +1.48 & 4 / 8 / 24 \\
Corruption & Fog & 36 & +0.39 & -0.56 to +3.10 & 6 / 0 / 30 \\
Corruption & JPEG & 36 & -0.20 & -4.23 to +2.01 & 6 / 7 / 23 \\
Severity & 1 & 48 & +0.28 & -1.03 to +2.82 & 8 / 2 / 38 \\
Severity & 3 & 48 & +0.16 & -3.37 to +2.09 & 10 / 8 / 30 \\
Severity & 5 & 48 & -0.51 & -4.28 to +3.10 & 7 / 16 / 25 \\
\bottomrule
\end{tabular}
}
\end{table}

\FloatBarrier

\section{Relative-retention scale sensitivity}

Table~\ref{tab:relative-retention} compares the AP-point interaction $\DeltaE$
with the relative-retention interaction $\DeltaR$. The exploratory balanced
means are $-0.02$ AP points and $+0.75$ pp, respectively, and 18 of 144 cell
point estimates differ in sign. The untouched holdouts have means of $-0.55$
AP points and $-0.12$ pp, with seven of 72 point-sign disagreements. These
differences are expected because the ratio estimand normalizes each treatment
by a different matched-clean denominator; they reinforce reporting absolute
corrupted AP and $\DeltaE$ rather than substituting one normalized score.

\begin{table}[htbp]
\centering
\caption{Relative-retention sensitivity. Mean $\DeltaE$ is in AP points and
mean $\DeltaR$ is in percentage points (pp). Sign differences compare finite
cell point estimates, not interval decisions.}
\label{tab:relative-retention}
\small
\resizebox{0.92\textwidth}{!}{% Generated by analysis/build_cviu_revision_artifacts.py.
\begin{tabular}{lrrrr}
\toprule
Scope & Cells & Mean $\Delta E$ (AP) & Mean $\Delta R$ (pp) & Sign differences \\
\midrule
Exploratory grid & 144 & -0.02 & +0.75 & 18 \\
Exploratory: COCO & 36 & -0.19 & +0.16 & 1 \\
Exploratory: VOC & 36 & -0.04 & +0.35 & 6 \\
Exploratory: KITTI & 36 & -0.46 & +0.49 & 6 \\
Exploratory: TT100K & 36 & +0.61 & +1.98 & 5 \\
Final holdouts & 72 & -0.55 & -0.12 & 7 \\
Final: VOC & 36 & -0.43 & -0.34 & 3 \\
Final: KITTI & 36 & -0.67 & +0.10 & 4 \\
\bottomrule
\end{tabular}
}
\end{table}

\begin{table}[htbp]
\centering
\caption{All 18 exploratory cells for which the point estimates of $\DeltaE$
and $\DeltaR$ have different signs. The seven holdout disagreements are
summarized in Table~\ref{tab:relative-retention}.}
\label{tab:relative-retention-disagreements}
\scriptsize
\resizebox{\textwidth}{!}{% Generated by analysis/build_cviu_revision_artifacts.py.
\begin{tabular}{lllrrr}
\toprule
Dataset & Model & Corruption & Severity & $\Delta E$ (AP) & $\Delta R$ (pp) \\
\midrule
COCO & YOLO11n & Gaussian noise & 5 & -0.31 & +1.23 \\
VOC & YOLO11m & Fog & 3 & -0.02 & +0.01 \\
VOC & YOLO11m & Gaussian noise & 3 & -0.23 & +0.14 \\
VOC & YOLO11m & Motion blur & 3 & -0.34 & +1.13 \\
VOC & YOLO11m & Motion blur & 5 & -1.09 & +0.67 \\
VOC & YOLO11n & Gaussian noise & 5 & -0.11 & +0.66 \\
VOC & YOLO11n & JPEG & 5 & -0.10 & +0.17 \\
KITTI & YOLO11m & Gaussian noise & 3 & -0.33 & +3.01 \\
KITTI & YOLO11m & Gaussian noise & 5 & -3.66 & +0.23 \\
KITTI & YOLO11m & JPEG & 5 & -0.16 & +0.06 \\
KITTI & YOLO11n & Gaussian noise & 5 & -1.13 & +1.21 \\
KITTI & YOLO11n & Motion blur & 5 & -1.70 & +0.36 \\
KITTI & YOLO11x & Motion blur & 5 & -1.18 & +0.15 \\
TT100K & YOLO11m & Gaussian noise & 5 & -1.21 & +0.03 \\
TT100K & YOLO11m & Motion blur & 5 & -0.26 & +1.95 \\
TT100K & YOLO11n & Gaussian noise & 5 & -0.72 & +0.57 \\
TT100K & YOLO11n & Motion blur & 5 & -0.03 & +3.78 \\
TT100K & YOLO11x & Motion blur & 5 & -0.03 & +0.07 \\
\bottomrule
\end{tabular}
}
\end{table}

\FloatBarrier

\section{Size-specific absolute-accuracy guardrail}

Table~\ref{tab:size-guardrail} binds relative size interactions to absolute
matched-clean AP, corrupted AP, and clean-to-corrupted loss. The AP$<5$
inventory makes explicit when an apparent contraction occurs close to a task
failure floor. COCO, VOC, and KITTI use COCO-style area endpoints; TT100K uses
its original-box-height endpoints.

\begin{table}[htbp]
\centering
\caption{Size-specific absolute-accuracy guardrail. Values are balanced means
across 36 model--corruption--severity cells per row.}
\label{tab:size-guardrail}
\scriptsize
\resizebox{\textwidth}{!}{\begin{tabular}{lllrrrrr}
\toprule
Dataset & Format & Endpoint & Cells & Clean AP & Corrupted AP & Mean $D$ (AP pt) & AP<5 \\
\midrule
COCO & INT8 & coco-area:small & 36 & 28.05 & 18.52 & +9.53 & 5 \\
COCO & INT8 & coco-area:large & 36 & 61.75 & 49.36 & +12.39 & 0 \\
COCO & FP8 & coco-area:small & 36 & 29.01 & 19.37 & +9.64 & 4 \\
COCO & FP8 & coco-area:large & 36 & 63.33 & 50.94 & +12.38 & 0 \\
KITTI & INT8 & coco-area:small & 36 & 50.90 & 34.18 & +16.72 & 6 \\
KITTI & INT8 & coco-area:large & 36 & 74.89 & 54.12 & +20.77 & 2 \\
KITTI & FP8 & coco-area:small & 36 & 54.02 & 36.52 & +17.50 & 6 \\
KITTI & FP8 & coco-area:large & 36 & 76.25 & 55.64 & +20.61 & 2 \\
TT100K & INT8 & original-height:XS & 36 & 12.02 & 6.42 & +5.60 & 20 \\
TT100K & INT8 & original-height:S & 36 & 38.03 & 25.46 & +12.57 & 8 \\
TT100K & INT8 & original-height:L & 36 & 58.50 & 43.67 & +14.83 & 1 \\
TT100K & INT8 & original-height:XL & 36 & 65.22 & 51.82 & +13.40 & 0 \\
TT100K & FP8 & original-height:XS & 36 & 14.00 & 6.93 & +7.06 & 21 \\
TT100K & FP8 & original-height:S & 36 & 40.15 & 27.11 & +13.04 & 8 \\
TT100K & FP8 & original-height:L & 36 & 58.52 & 44.97 & +13.55 & 0 \\
TT100K & FP8 & original-height:XL & 36 & 66.42 & 53.50 & +12.93 & 0 \\
VOC & INT8 & coco-area:small & 36 & 28.12 & 19.98 & +8.14 & 6 \\
VOC & INT8 & coco-area:large & 36 & 71.96 & 54.99 & +16.97 & 0 \\
VOC & FP8 & coco-area:small & 36 & 28.97 & 20.51 & +8.46 & 6 \\
VOC & FP8 & coco-area:large & 36 & 72.97 & 56.12 & +16.85 & 0 \\
\bottomrule
\end{tabular}
}
\end{table}

For TT100K, Table~\ref{tab:tt100k-combined-guardrail} reports the combined
small-like and large-like endpoints used by the height interaction itself.
These endpoints are evaluated directly; they are not arithmetic averages of
the component XS/S or L/XL AP values, and the medium-height band is excluded by
the frozen endpoint definition.

\begin{table}[htbp]
\centering
\caption{TT100K absolute-accuracy guardrail for the combined original-height
endpoints used in $\DeltaPsi$. $D$ is the signed clean-to-corrupted AP loss.}
\label{tab:tt100k-combined-guardrail}
\small
\resizebox{0.94\textwidth}{!}{% Generated by analysis/build_cviu_revision_artifacts.py.
% 36-record TT100K paired-contrast set SHA-256: d318a245fcd380ec8e11d3beedc4eec93c2e66ae30a182a37e9312df34f338e6
% Three JPEG-95 clean-arm-cache set SHA-256: 9db85b0ac72c077b11d4848bea90e293951763342112a5386ce48d2910d04af5
% Annotation SHA-256: aedb7cccb2135dd455ce8ee3a4b1feeb019670f029fee8cba5577b390e579a80
% Means were audited with the same maxDets=100 evaluator used for Delta Psi; medium [24,48) is excluded.
% Combined endpoints were evaluated directly and are not arithmetic averages of XS/S or L/XL AP.
\begin{tabular}{llrrrrrr}
\toprule
Treatment & Original-height endpoint & Cells & Clean AP & Corrupted AP & Mean $D$ (AP) & AP$<5$ & AP$<10$ \\
\midrule
INT8 & small-like [0,24) & 36 & 37.74 & 25.25 & +12.49 & 8 & 9 \\
INT8 & large-like [48,$\infty$) & 36 & 55.61 & 41.57 & +14.04 & 1 & 2 \\
FP8 & small-like [0,24) & 36 & 39.92 & 26.91 & +13.01 & 8 & 9 \\
FP8 & large-like [48,$\infty$) & 36 & 56.29 & 43.11 & +13.18 & 0 & 2 \\
\bottomrule
\end{tabular}
}
\end{table}

\FloatBarrier

\section{Codec-control sensitivity}

The primary estimand uses deterministic JPEG-95 matched-clean inputs so that
both clean and corrupted arms share the terminal materialization codec.
Table~\ref{tab:codec} compares original-source clean AP with JPEG-95 clean AP.
This point sensitivity quantifies the control choice; it is not a bootstrap
interval and does not replace the codec-matched estimand.

\begin{table}[htbp]
\centering
\caption{Original-source clean versus deterministic JPEG-95 clean AP. The last
column is JPEG-95 minus original-source AP, all on the 0--100 AP-point scale.}
\label{tab:codec}
\scriptsize
\resizebox{\textwidth}{!}{\begin{tabular}{lllrrr}
\toprule
Dataset & Model & Precision & Original clean AP (points) & JPEG-95 clean AP (points) & Codec minus original (AP points) \\
\midrule
COCO & YOLO11n & FP32 & 38.73 & 38.70 & -0.035 \\
COCO & YOLO11n & INT8 & 36.64 & 36.65 & +0.014 \\
COCO & YOLO11n & FP8 & 37.74 & 37.81 & +0.072 \\
COCO & YOLO11m & FP32 & 50.68 & 50.57 & -0.113 \\
COCO & YOLO11m & INT8 & 48.85 & 48.82 & -0.025 \\
COCO & YOLO11m & FP8 & 50.08 & 50.05 & -0.028 \\
COCO & YOLO11x & FP32 & 53.81 & 53.69 & -0.118 \\
COCO & YOLO11x & INT8 & 52.75 & 52.56 & -0.197 \\
COCO & YOLO11x & FP8 & 53.52 & 53.47 & -0.056 \\
VOC & YOLO11n & FP32 & 62.54 & 62.56 & +0.013 \\
VOC & YOLO11n & INT8 & 61.42 & 61.51 & +0.090 \\
VOC & YOLO11n & FP8 & 62.06 & 62.19 & +0.125 \\
VOC & YOLO11m & FP32 & 68.53 & 68.56 & +0.036 \\
VOC & YOLO11m & INT8 & 66.75 & 66.66 & -0.089 \\
VOC & YOLO11m & FP8 & 68.31 & 68.36 & +0.044 \\
VOC & YOLO11x & FP32 & 70.09 & 70.05 & -0.032 \\
VOC & YOLO11x & INT8 & 69.07 & 69.03 & -0.047 \\
VOC & YOLO11x & FP8 & 70.00 & 69.94 & -0.060 \\
KITTI & YOLO11n & FP32 & 64.79 & 64.84 & +0.049 \\
KITTI & YOLO11n & INT8 & 61.52 & 61.69 & +0.170 \\
KITTI & YOLO11n & FP8 & 63.68 & 63.83 & +0.147 \\
KITTI & YOLO11m & FP32 & 73.82 & 73.87 & +0.049 \\
KITTI & YOLO11m & INT8 & 69.22 & 69.34 & +0.124 \\
KITTI & YOLO11m & FP8 & 73.32 & 73.31 & -0.016 \\
KITTI & YOLO11x & FP32 & 72.93 & 73.06 & +0.127 \\
KITTI & YOLO11x & INT8 & 70.87 & 70.82 & -0.053 \\
KITTI & YOLO11x & FP8 & 72.20 & 72.26 & +0.061 \\
TT100K & YOLO11n & FP32 & 28.78 & 28.76 & -0.013 \\
TT100K & YOLO11n & INT8 & 25.10 & 25.29 & +0.191 \\
TT100K & YOLO11n & FP8 & 27.12 & 26.41 & -0.709 \\
TT100K & YOLO11m & FP32 & 56.57 & 56.55 & -0.014 \\
TT100K & YOLO11m & INT8 & 53.05 & 52.97 & -0.074 \\
TT100K & YOLO11m & FP8 & 54.75 & 54.40 & -0.350 \\
TT100K & YOLO11x & FP32 & 57.63 & 57.66 & +0.034 \\
TT100K & YOLO11x & INT8 & 55.50 & 56.02 & +0.521 \\
TT100K & YOLO11x & FP8 & 56.76 & 56.10 & -0.667 \\
\bottomrule
\end{tabular}
}
\end{table}

\FloatBarrier

\section{Weighting, omission, and runtime sensitivities}

Alternative weighting changes the descriptive target, not merely the standard
error. Table~\ref{tab:weighting} therefore reports equal-cell, image-count, and
object-count means together with leave-one-dataset and leave-one-corruption
summaries. Sign changes near zero expose conditional heterogeneity; they do not
identify a causal role for the omitted factor.

\begin{table}[htbp]
\centering
\caption{Weighting and leave-one-factor-out sensitivity of the 144-cell ledger.}
\label{tab:weighting}
\scriptsize
\resizebox{0.92\textwidth}{!}{\begin{tabular}{llrrl}
\toprule
Analysis & Level & Cells & Mean $\Delta E$ (AP pt) & SD / IQR (AP pt) \\
\midrule
Cell distribution & all cells & 144 & -0.02 & 1.20 / [-0.39, +0.56] \\
Weighting & equal cell & 144 & -0.02 & -- \\
Weighting & image weighted & 144 & +0.00 & -- \\
Weighting & object weighted & 144 & -0.09 & -- \\
Leave-one-dataset-out & COCO & 108 & +0.03 & -- \\
Leave-one-dataset-out & KITTI & 108 & +0.13 & -- \\
Leave-one-dataset-out & TT100K & 108 & -0.23 & -- \\
Leave-one-dataset-out & VOC & 108 & -0.01 & -- \\
Leave-one-corruption-out & Fog & 108 & -0.16 & -- \\
Leave-one-corruption-out & Gaussian noise & 108 & -0.02 & -- \\
Leave-one-corruption-out & JPEG & 108 & +0.04 & -- \\
Leave-one-corruption-out & Motion blur & 108 & +0.05 & -- \\
\bottomrule
\end{tabular}
}
\end{table}

Runtime is reported only for matched engine conditions. Ratios exclude image
decoding, preprocessing, host--device transfer, detector decoding, NMS, power,
and energy. Table~\ref{tab:runtime-sensitivity} stratifies the engine-only
measurements by input size and YOLO11 capacity rather than pooling unlike
workloads into one deployment claim.

\begin{table}[htbp]
\centering
\caption{Matched engine-only latency-ratio sensitivity. Entries are median
[Q1, Q3] across retained conditions.}
\label{tab:runtime-sensitivity}
\small
\resizebox{0.86\textwidth}{!}{\begin{tabular}{llrlll}
\toprule
Stratum & Level & Conditions & FP32/INT8 & FP32/FP8 & INT8/FP8 \\
\midrule
Input size & 640 & 9 & 3.08 [1.56, 3.08] & 3.32 [1.81, 3.33] & 1.09 [1.08, 1.16] \\
Input size & 1280 & 3 & 2.81 [2.38, 2.84] & 2.86 [2.33, 2.91] & 1.02 [0.97, 1.02] \\
Capacity & YOLO11m & 4 & 3.08 [3.01, 3.09] & 3.34 [3.21, 3.35] & 1.08 [1.06, 1.09] \\
Capacity & YOLO11n & 4 & 1.55 [1.52, 1.65] & 1.80 [1.79, 1.81] & 1.17 [1.10, 1.18] \\
Capacity & YOLO11x & 4 & 3.06 [3.00, 3.09] & 3.32 [3.22, 3.33] & 1.07 [1.05, 1.08] \\
\bottomrule
\end{tabular}
}
\end{table}

Table~\ref{tab:runtime-conditions} exposes all 12 matched conditions underlying
the ratio summaries. Ratios are computed condition by condition; their medians
therefore need not equal ratios formed from separately pooled marginal medians.

\begin{table}[htbp]
\centering
\caption{Condition-level TensorRT engine-only latency and matched ratios for
the primary YOLO11 grid.}
\label{tab:runtime-conditions}
\scriptsize
\resizebox{\textwidth}{!}{% Generated by analysis/build_cviu_revision_artifacts.py.
\begin{tabular}{llrrrrrr}
\toprule
Dataset & Model & FP32 ms & INT8 ms & FP8 ms & 32/INT8 & 32/FP8 & INT8/FP8 \\
\midrule
COCO & YOLO11m & 2.445 & 0.792 & 0.734 & 3.09 & 3.33 & 1.08 \\
COCO & YOLO11n & 0.543 & 0.365 & 0.306 & 1.49 & 1.78 & 1.19 \\
COCO & YOLO11x & 5.336 & 1.753 & 1.613 & 3.04 & 3.31 & 1.09 \\
VOC & YOLO11m & 2.422 & 0.786 & 0.725 & 3.08 & 3.34 & 1.08 \\
VOC & YOLO11n & 0.550 & 0.354 & 0.305 & 1.56 & 1.80 & 1.16 \\
VOC & YOLO11x & 5.367 & 1.709 & 1.609 & 3.14 & 3.33 & 1.06 \\
KITTI & YOLO11m & 2.422 & 0.785 & 0.720 & 3.08 & 3.36 & 1.09 \\
KITTI & YOLO11n & 0.551 & 0.358 & 0.304 & 1.54 & 1.81 & 1.18 \\
KITTI & YOLO11x & 5.322 & 1.730 & 1.602 & 3.08 & 3.32 & 1.08 \\
TT100K & YOLO11m & 5.537 & 1.970 & 1.936 & 2.81 & 2.86 & 1.02 \\
TT100K & YOLO11n & 1.406 & 0.723 & 0.785 & 1.95 & 1.79 & 0.92 \\
TT100K & YOLO11x & 12.965 & 4.516 & 4.393 & 2.87 & 2.95 & 1.03 \\
\bottomrule
\end{tabular}
}
\end{table}

\FloatBarrier

\section{Selection-disjoint holdout and corruption-realization sensitivities}

The selection-disjoint branch froze disjoint checkpoint-selection and final image lists
before training. VOC used 8,218 training, 2,510 selection, and 5,823 final
images. For KITTI, the 1,496-image selection partition was reserved first; a
deterministic class-covered 20\% subset of the remaining 5,985 images formed
the 1,197-image final set, and the other 4,788 images formed the training set
(Table~\ref{tab:holdout-splits}). Calibration identities were drawn only
from training data. The final factorial rerun retained 72 equal-weight
VOC/KITTI cells and reused a common 2,000-draw schedule within each dataset.
Its balanced direct interaction was -0.55 AP points with percentile interval
$[-0.78,-0.29]$; VOC was -0.43 $[-0.54,-0.29]$ and KITTI was -0.67
$[-1.13,-0.20]$. FP8 retained 1.60 more matched-clean AP on average, and the
clean advantage was positive in all six dataset--capacity blocks. Nine of 144
corrupted format arms were below 5 AP and 18 were below 10 AP. This branch
removes evaluation-after-selection reuse for these two datasets, but it does
not convert the broader exploratory grid into a prospective confirmatory study.

\begin{table}[htbp]
\centering
\caption{Selection-disjoint partition counts. The KITTI final set is a
class-covered subset of the post-selection remainder, not 20\% of the complete
7,481-image labeled universe.}
\label{tab:holdout-splits}
\small
\resizebox{0.88\textwidth}{!}{% Generated by analysis/build_cviu_revision_artifacts.py.
\begin{tabular}{llrrr}
\toprule
Dataset & Partition & Images & Objects & Observed/declared classes \\
\midrule
VOC & Train & 8,218 & 19,910 & 20/20 \\
VOC & Selection & 2,510 & 6,307 & 20/20 \\
VOC & Final & 5,823 & 13,841 & 20/20 \\
KITTI & Train & 4,788 & 25,847 & 8/8 \\
KITTI & Selection & 1,496 & 8,128 & 8/8 \\
KITTI & Final & 1,197 & 6,595 & 8/8 \\
\bottomrule
\end{tabular}
}
\end{table}

The frozen reference ladder was checked on the final partitions before the
quantized comparisons. Table~\ref{tab:holdout-parity} reports absolute AP-point
gaps; all six blocks satisfied the recorded parity gate.

\begin{table}[htbp]
\centering
\caption{Selection-disjoint reference-ladder parity. Values are absolute
AP-point gaps between the named implementations.}
\label{tab:holdout-parity}
\small
\resizebox{0.9\textwidth}{!}{% Generated by analysis/build_cviu_revision_artifacts.py; absolute AP-point gaps.
\begin{tabular}{llrrrr}
\toprule
Dataset & Model & Source--ORT & Source--TRT32 & TRT32--TRT16 & Gate \\
\midrule
VOC & YOLO11m & 0.0013 & 0.0024 & 0.0004 & Pass \\
VOC & YOLO11n & 0.0027 & 0.0052 & 0.0008 & Pass \\
VOC & YOLO11x & 0.0041 & 0.0042 & 0.0213 & Pass \\
KITTI & YOLO11m & 0.0362 & 0.0404 & 0.1084 & Pass \\
KITTI & YOLO11n & 0.0221 & 0.0036 & 0.0123 & Pass \\
KITTI & YOLO11x & 0.0001 & 0.0126 & 0.2554 & Pass \\
\bottomrule
\end{tabular}
}
\end{table}

The corruption-realization sensitivity used three fixed realization seeds over
36 base conditions on class-covering 512-image subsets of VOC, KITTI, and
TT100K. Its equal-cell mean $\DeltaE$ was -0.23 AP points, with nested
realization/image percentiles $[-0.95,-0.28,+0.37]$. Mean between-realization
variance pooled over stochastic and deterministic conditions was 0.3989
AP-points$^2$, compared with 0.7446 within-image-bootstrap variance. The more
appropriate stochastic-only between-realization mean was 0.5318
AP-points$^2$; deterministic JPEG has no realization variance by construction
and is marked not applicable in Table~\ref{tab:realization-variance}. For
$\DeltaPsi$, the mean was -1.04 with nested percentiles
$[-3.05,-1.42,+0.10]$. Only three materializations were used, so these are
conditional sensitivity summaries rather than population intervals over all
possible corruptions. In this separate sensitivity branch, stochastic fields
and motion-blur angles were nested across severity; in the exploratory grid,
severity remains an ordinal factor conditional on its recorded materialization.

\begin{table}[htbp]
\centering
\caption{Realization-specific direct interactions. SD is the descriptive
standard deviation across the three fixed realization labels.}
\label{tab:realization-seeds}
\small
\resizebox{0.82\textwidth}{!}{% Generated by analysis/build_cviu_revision_artifacts.py.
\begin{tabular}{lrrr}
\toprule
Realization seed & Cells & Mean $\Delta E$ (AP) & Mean $\Delta\Psi$ (AP) \\
\midrule
202608181 & 36 & -0.12 & -1.27 \\
202608182 & 36 & -0.26 & -0.83 \\
202608183 & 36 & -0.32 & -1.01 \\
Across-seed summary & 108 & -0.23; SD 0.10 & -1.04; SD 0.22 \\
\bottomrule
\end{tabular}
}
\end{table}

\begin{table}[htbp]
\centering
\caption{Per-corruption realization sensitivity. Variances are in squared AP
points. JPEG is deterministic and is excluded from the stochastic-only mean.}
\label{tab:realization-variance}
\small
\resizebox{0.92\textwidth}{!}{% Generated by analysis/build_cviu_revision_artifacts.py; variances use AP-points squared.
\begin{tabular}{lrrr}
\toprule
Corruption & Mean $\Delta E$ (AP) & Between-realization variance & Within-image variance \\
\midrule
Gaussian noise & -0.12 & 0.5960 & 0.7306 \\
Motion blur & -0.86 & 0.5931 & 0.8302 \\
Fog & +0.23 & 0.4063 & 0.7114 \\
JPEG & -0.18 & n/a (deterministic) & 0.7062 \\
Stochastic-only mean & -- & 0.5318 & -- \\
\bottomrule
\end{tabular}
}
\end{table}

\begin{table}[htbp]
\centering
\caption{Scope and uncertainty target of the four conditional evidence layers.
The reported intervals are not interchangeable estimates of one parameter.}
\label{tab:conditionality-scope}
\scriptsize
\resizebox{\textwidth}{!}{% Generated by analysis/build_cviu_revision_artifacts.py. Intervals have different conditional scopes.
\begin{tabular}{llrrll}
\toprule
Analysis & Dataset/model/severity scope & Cells & Mean $\Delta E$ & Interval (AP) & Uncertainty target \\
\midrule
Exploratory grid & 4; n/m/x; S1/3/5 & 144 & -0.02 & [-0.24, +0.15] & Paired image bootstrap \\
Final holdouts & VOC/KITTI; n/m/x; S1/3/5 & 72 & -0.55 & [-0.78, -0.29] & Paired image bootstrap \\
Realization subset & 3 datasets; m; S1/3/5 & 36 base / 108 & -0.23 & [-0.95, +0.37] & 3-label nested sensitivity \\
Training/calibration seeds & 3 datasets; m; S5 & 108 & -1.11 & not estimated & Descriptive $3\times3$ grid \\
\bottomrule
\end{tabular}
}
\end{table}

\FloatBarrier

\section{Training--calibration seed sensitivity}

The targeted $3\times3$ training--calibration grid covers YOLO11m on VOC,
KITTI, and TT100K at severity 5. Across 108 cells, the equally weighted mean
$\DeltaE$ was -1.11 AP points. Training-seed marginal means ranged from -1.64
to -0.68, and calibration-seed marginal means from -1.32 to -0.84. The
original-seed intersection was -1.21, close to the balanced seed-grid mean.
These results do not extend to COCO, other capacity rungs, lower severities, or
an unspecified population of training and calibration seeds.

Table~\ref{tab:training-realizations} records the optimizer realization and the
checkpoint-selection outcome for each exploratory transfer or untouched-holdout
training run for which training was performed in this study. COCO uses an
official pretrained checkpoint and is therefore absent. Selection AP is not a
final-evaluation result.

\begin{table}[htbp]
\centering
\caption{Training realizations and selected best epochs. Effective weight
decay records the realized batch/accumulation scaling. AP values are on the
0--100 scale and refer only to the checkpoint-selection partition.}
\label{tab:training-realizations}
\scriptsize
\resizebox{\textwidth}{!}{% Generated by analysis/build_cviu_revision_artifacts.py.
% Exploratory manifests set SHA-256: f5ffaff0c5416b88e3e425f7ea126558f8ed2ddb508fdf754ffa2b6e6cb70806
% Exploratory results.csv set SHA-256: 3ef86140dc8f60017c944541accc964bd88b38135278447cd2f4791f40f48a95
% Exploratory training-log set SHA-256: fbd98d722dfccce1e404d0f6de38cb410126aa648e9e60782c704378d836c077
% Held-out training-manifest set SHA-256: a32c00a867d4923226c7c9ff6fa187b19528cb3e5ea2bb2a48913af4fd1162eb
% Held-out best-only-manifest set SHA-256: d70a474dfa8569ba52b2f77bb59531077860ff2f6ebc9a7c6cad55bc640494d1
% Held-out best epochs/AP are a fixed transcription of execution-host results.csv files audited on 2026-09-03;
% the compact workspace binds args.yaml/best.pt hashes but does not retain those six results.csv files.
\begin{tabular}{lllrlrrrrr}
\toprule
Branch & Dataset & Model & Batch & Optimizer & Base LR & Mom./$\beta_1$ & Effective WD & Best epoch & Selection AP$_{50:95}$ \\
\midrule
Expl. & VOC & YOLO11n & 52 & MuSGD & 0.01 & 0.9 & 0.00040625 & 99 & 63.653 \\
Expl. & VOC & YOLO11m & 14 & MuSGD & 0.01 & 0.9 & 0.000546875 & 65 & 69.593 \\
Expl. & VOC & YOLO11x & 6 & MuSGD & 0.01 & 0.9 & 0.000515625 & 67 & 71.228 \\
Expl. & KITTI & YOLO11n & 55 & AdamW & 0.000833 & 0.9 & 0.0004296875 & 100 & 65.203 \\
Expl. & KITTI & YOLO11m & 14 & AdamW & 0.000833 & 0.9 & 0.000546875 & 91 & 73.850 \\
Expl. & KITTI & YOLO11x & 6 & AdamW & 0.000833 & 0.9 & 0.000515625 & 100 & 73.012 \\
Expl. & TT100K & YOLO11n & 11 & AdamW & 4.4e-05 & 0.9 & 0.000515625 & 81 & 7.197 \\
Expl. & TT100K & YOLO11m & 3 & AdamW & 4.4e-05 & 0.9 & 0.0004921875 & 74 & 26.746 \\
Expl. & TT100K & YOLO11x & 1 & AdamW & 4.4e-05 & 0.9 & 0.0005 & 78 & 26.117 \\
Holdout & VOC & YOLO11n & 16 & MuSGD & 0.01 & 0.9 & 0.0005 & 82 & 57.925 \\
Holdout & VOC & YOLO11m & 16 & MuSGD & 0.01 & 0.9 & 0.0005 & 64 & 63.428 \\
Holdout & VOC & YOLO11x & 16 & MuSGD & 0.01 & 0.9 & 0.0005 & 73 & 63.662 \\
Holdout & KITTI & YOLO11n & 16 & AdamW & 0.000833 & 0.9 & 0.0005 & 100 & 62.026 \\
Holdout & KITTI & YOLO11m & 16 & AdamW & 0.000833 & 0.9 & 0.0005 & 99 & 71.211 \\
Holdout & KITTI & YOLO11x & 16 & AdamW & 0.000833 & 0.9 & 0.0005 & 98 & 71.811 \\
\bottomrule
\end{tabular}
}
\end{table}

\begin{table}[htbp]
\centering
\caption{Descriptive two-way seed decomposition within each targeted
YOLO11m severity-5 block. The last three columns are percentages of the
within-block total sum of squares; Train$\times$cal. denotes non-additivity.}
\label{tab:multiseed-decomposition}
\scriptsize
\resizebox{0.86\textwidth}{!}{%
\begin{tabular}{llrrrr}
\toprule
Dataset & Corruption & Mean & Train (\%) & Cal. (\%) & Train$\times$cal. (\%) \\
\midrule
VOC & Gaussian noise & -1.494 & 28.6 & 23.0 & 48.4 \\
VOC & Motion blur & -1.256 & 11.3 & 18.0 & 70.7 \\
VOC & Fog & +0.018 & 34.2 & 36.5 & 29.3 \\
VOC & JPEG & -2.033 & 9.3 & 80.2 & 10.5 \\
KITTI & Gaussian noise & -3.532 & 51.5 & 13.8 & 34.8 \\
KITTI & Motion blur & -4.350 & 57.1 & 37.4 & 5.5 \\
KITTI & Fog & +1.608 & 23.9 & 24.8 & 51.3 \\
KITTI & JPEG & +0.120 & 38.4 & 54.6 & 7.0 \\
TT100K & Gaussian noise & -1.381 & 8.7 & 22.5 & 68.8 \\
TT100K & Motion blur & -0.998 & 12.6 & 5.5 & 82.0 \\
TT100K & Fog & +0.769 & 11.5 & 27.5 & 61.0 \\
TT100K & JPEG & -0.787 & 90.9 & 0.5 & 8.5 \\
\bottomrule
\end{tabular}}
\end{table}

\FloatBarrier

\section{Fixed-class-universe bootstrap sensitivity}

For the targeted TT100K YOLO11m/motion-blur/severity-5 stress cell, a
positive-weight Bayesian image bootstrap retained the complete category
universe across all 3,067 images. Each draw assigned independent unit-rate
exponential image weights, normalized the 3,067 weights to mean one, applied
them to ground-truth positives and true/false-positive contributions, and
reused the same vector across all four arms. With 10,000 draws, the $\DeltaE$ percentiles
were $[-1.19,-0.17,+0.88]$, compared with $[-1.79,-0.20,+1.34]$ under the
ordinary 2,000-draw image bootstrap. The height-based $\DeltaPsi$ percentiles
were $[-5.69,-4.08,-2.48]$ versus $[-7.22,-4.48,-1.81]$. The median sign and
zero-crossing status of $\DeltaE$, and the negative-interval status of
$\DeltaPsi$, were unchanged. This targeted cell does not establish
class-composition invariance for the full TT100K grid.

\FloatBarrier

\section{Architecture-portability stress cases}

The RT-DETR-L and RetinaNet-R50-FPN-v2 extension tests whether the frozen
quantization recipes transfer beyond the YOLO family. It is interpreted in a
fixed sequence: matched-clean fidelity first, absolute corrupted AP second,
and the corruption interaction third.

\begin{table}[htbp]
\centering
\caption{Matched JPEG-95 clean AP and clean quantization gaps for the
architecture-portability stress cases.}
\label{tab:cross-clean}
\scriptsize
\resizebox{\textwidth}{!}{\begin{tabular}{llrrrrr}
\toprule
Dataset & Family & FP32 AP & INT8 AP & FP8 AP & $Q_{95}^{\mathrm{INT8}}$ & $Q_{95}^{\mathrm{FP8}}$ \\
\midrule
VOC & RT-DETR-L & 70.44 & 1.14 & 70.14 & 69.30 & 0.30 \\
VOC & RetinaNet-R50-FPN-v2 & 57.22 & 33.69 & 56.17 & 23.53 & 1.05 \\
KITTI & RT-DETR-L & 12.29 & 3.32 & 12.02 & 8.97 & 0.27 \\
KITTI & RetinaNet-R50-FPN-v2 & 50.43 & 7.13 & 49.27 & 43.30 & 1.15 \\
TT100K & RT-DETR-L & 43.17 & 5.46 & 42.13 & 37.71 & 1.04 \\
TT100K & RetinaNet-R50-FPN-v2 & 3.37 & 2.60 & 3.38 & 0.77 & -0.01 \\
\bottomrule
\end{tabular}}
\end{table}

\begin{table}[htbp]
\centering
\caption{Canonical cross-family interaction summaries under
$E=Q_c-Q_{\Jclean}$. Negative INT8 means largely reflect contraction of a large
clean gap near low corrupted AP.}
\label{tab:cross-interaction}
\scriptsize
\resizebox{\textwidth}{!}{% Generated by scripts/rebuild_cross_family_evidence.py; do not edit manually.
\begin{tabular}{llrrrrr}
\toprule
Family & Format & Mean $Q_{95}$ & Mean $E$ & $E$ range & Mean corrupted AP & AP$<5$ \\
\midrule
RT-DETR-L & INT8 & 38.66 & -6.81 & -54.39--+1.22 & 1.97 & 34/36 \\
RT-DETR-L & FP8 & 0.54 & -0.02 & -1.06--+2.01 & 33.30 & 4/36 \\
RetinaNet-R50-FPN-v2 & INT8 & 22.53 & -6.21 & -38.86--+1.14 & 11.34 & 16/36 \\
RetinaNet-R50-FPN-v2 & FP8 & 0.73 & +0.01 & -1.11--+0.64 & 26.91 & 12/36 \\
\bottomrule
\end{tabular}
}
\end{table}

\begin{table}[htbp]
\centering
\caption{Descriptive direct $\DeltaE$ summaries for the architecture stress
cases. These values are not evidence of intrinsic datatype superiority.}
\label{tab:cross-direct}
\small
\resizebox{0.72\textwidth}{!}{% Generated by scripts/rebuild_cross_family_evidence.py; do not edit manually.
\begin{tabular}{lrrrr}
\toprule
Family & Cells & Mean $\Delta E$ & Minimum & Maximum \\
\midrule
RT-DETR-L & 36 & -6.79 & -54.33 & +1.51 \\
RetinaNet-R50-FPN-v2 & 36 & -6.23 & -37.75 & +1.16 \\
\bottomrule
\end{tabular}
}
\end{table}

Direct contrasts are computed from unrounded cell values; consequently,
displayed rounded $E_{\mathrm{INT8}}$ and $E_{\mathrm{FP8}}$ components need not
subtract to the displayed direct $\DeltaE$ exactly.

\begin{table}[htbp]
\centering
\caption{Per-dataset matched engine-only runtime for the architecture stress
cases. Ratios are formed within each dataset and family.}
\label{tab:cross-runtime-dataset}
\scriptsize
\resizebox{\textwidth}{!}{% Generated by analysis/build_cviu_revision_artifacts.py.
\begin{tabular}{llrrrrr}
\toprule
Family & Dataset & FP32 ms & INT8 ms & FP8 ms & 32/INT8 & 32/FP8 \\
\midrule
RT-DETR-L & VOC & 2.770 & 1.120 & 1.203 & 2.47 & 2.30 \\
RT-DETR-L & KITTI & 2.756 & 1.121 & 1.197 & 2.46 & 2.30 \\
RT-DETR-L & TT100K & 6.814 & 2.434 & 2.591 & 2.80 & 2.63 \\
RetinaNet-R50-FPN-v2 & VOC & 5.080 & 1.113 & 1.290 & 4.56 & 3.94 \\
RetinaNet-R50-FPN-v2 & KITTI & 4.994 & 1.085 & 1.281 & 4.60 & 3.90 \\
RetinaNet-R50-FPN-v2 & TT100K & 14.285 & 3.613 & 4.232 & 3.95 & 3.38 \\
\bottomrule
\end{tabular}
}
\end{table}

FP8 remained within 1.15 matched-clean AP points of the recorded high-precision TensorRT reference in all
six dataset--family blocks. By contrast, the transferred INT8 recipe produced
mean clean gaps of 38.66 AP for RT-DETR-L and 22.53 for RetinaNet. Mean
corrupted INT8 AP was 1.97 and 11.34, respectively. The supported conclusion is
therefore a failure of recipe portability under these recorded treatments, not
an intrinsic limitation of INT8 or a universal advantage of FP8.

\FloatBarrier

\section{Evidence-generation and audit boundary}

The versioned public evidence archive cited in the main article separates three
layers:
\begin{itemize}
  \item \textbf{Metric ledgers and evidence reports} contain the retained
  summaries and upstream SHA-256 bindings supplied with the study.
  \item \textbf{Generation scripts} deterministically reconstruct the
  decision-impact diagnostic and the canonical cross-family sign convention
  from packaged CSV inputs.
  \item \textbf{LaTeX} consumes prevalidated tables and figures; compilation
  itself does not recompute model inference or independently verify external
  files that are absent from the compact journal source package.
\end{itemize}

The external research pipeline validated metric-to-run, input-manifest,
ordered-image, prediction, and analysis hashes before generating the packaged
artifacts. The manuscript does not claim that the LaTeX engine performs those
checks. The validator archived with the public research release checks
package-level consistency: row counts,
formula signs, required files, abstract/highlight constraints, graphical
abstract dimensions, unresolved placeholders, and SHA-256 manifest integrity.
Raw datasets, checkpoints, TensorRT engines, and per-image predictions are not
redistributed in this submission source.

The compact journal source is intentionally limited to manuscript sources,
figures, bibliography, and generated tables. It is not itself the evidence
archive and does not rerun archive validation during LaTeX compilation.

\FloatBarrier

\section{Minimum reporting checklist for paired corruption studies}

For reuse of this protocol, a report should state at least:
\begin{itemize}
  \item the exact clean materialization and whether corrupted arms share its
  terminal codec;
  \item the treatment-level quantization recipe, graph coverage, calibration
  identities, runtime version, and decoder;
  \item the atomic image IDs and proof that compared arms consume identical
  ordered inputs;
  \item the direct clean-adjusted contrast, not only corrupted treatment means;
  \item absolute corrupted AP and explicit low-accuracy floor inventories;
  \item the resampling unit, pairing structure, draw reuse, and uncertainty
  boundary;
  \item the finite-grid weighting target and all non-pooled endpoint families;
  \item sensitivity to split reuse, corruption realization, training and
  calibration seeds, class universe, and recipe transfer where feasible.
\end{itemize}

\FloatBarrier

\end{document}
